% !TEX TS-program = pdflatex
% !TEX encoding = UTF-8 Unicode

% This file is a template using the "beamer" package to create slides for a talk or presentation
% - Giving a talk on some subject.
% - The talk is between 15min and 45min long.
% - Style is ornate.

% MODIFIED by Jonathan Kew, 2008-07-06
% The header comments and encoding in this file were modified for inclusion with TeXworks.
% The content is otherwise unchanged from the original distributed with the beamer package.

\documentclass{beamer}


% Copyright 2004 by Till Tantau <tantau@users.sourceforge.net>.
%
% In principle, this file can be redistributed and/or modified under
% the terms of the GNU Public License, version 2.
%
% However, this file is supposed to be a template to be modified
% for your own needs. For this reason, if you use this file as a
% template and not specifically distribute it as part of a another
% package/program, I grant the extra permission to freely copy and
% modify this file as you see fit and even to delete this copyright
% notice. 


\mode<presentation>
{
  \usetheme{Warsaw}
  % or ...

  \setbeamercovered{transparent}
  % or whatever (possibly just delete it)
}


\usepackage[english]{babel}
% or whatever

\usepackage[utf8]{inputenc}
% or whatever

\usepackage{mathrsfs}

\usepackage{times}
\usepackage[T1]{fontenc}
% Or whatever. Note that the encoding and the font should match. If T1
% does not look nice, try deleting the line with the fontenc.

%%% MATH RELATED
\usepackage{amsmath,amsthm,amssymb}
\newtheorem{prob}{Problem}
\newtheorem{idea}{Idea}
\newtheorem{defn}{Definition}
\newtheorem{thm}{Theorem}
\newtheorem{conj}{Conjecture}
\newtheorem{ex}{Example}

\newcommand{\bbc}{\mathbb{C}}
\newcommand{\bbr}{\mathbb{R}}
\newcommand{\bbn}{\mathbb{N}}
\newcommand{\Ai}{\text{Ai}}
\newcommand{\Be}{\text{Be}}
\newcommand{\wt}[1]{\widetilde{#1}}

\title{Math GRE Prep Day 3} 
\subtitle
{} % (optional)

\author[W.R. Casper] % (optional, use only with lots of authors)
{W.R. Casper}
% - Use the \inst{?} command only if the authors have different
%   affiliation.

\institute[California State University Fullerton] % (optional, but mostly needed)
{
  Department of Mathematics\\
  California State University Fullerton}
% - Use the \inst command only if there are several affiliations.
% - Keep it simple, no one is interested in your street address.

\subject{Talks}
% This is only inserted into the PDF information catalog. Can be left
% out. 



% If you have a file called "university-logo-filename.xxx", where xxx
% is a graphic format that can be processed by latex or pdflatex,
% resp., then you can add a logo as follows:

% \pgfdeclareimage[height=0.5cm]{university-logo}{university-logo-filename}
% \logo{\pgfuseimage{university-logo}}



% Delete this, if you do not want the table of contents to pop up at
% the beginning of each subsection:
\AtBeginSubsection[]
{
  \begin{frame}<beamer>{Outline}
    \tableofcontents[currentsection,currentsubsection]
  \end{frame}
}


% If you wish to uncover everything in a step-wise fashion, uncomment
% the following command: 

%\beamerdefaultoverlayspecification{<+->}


\begin{document}

\begin{frame}
  \titlepage
\end{frame}

\begin{frame}{Outline}
  \tableofcontents
  % You might wish to add the option [pausesections]
\end{frame}

% Since this a solution template for a generic talk, very little can
% be said about how it should be structured. However, the talk length
% of between 15min and 45min and the theme suggest that you stick to
% the following rules:  

% - Exactly two or three sections (other than the summary).
% - At *most* three subsections per section.
% - Talk about 30s to 2min per frame. So there should be between about
%   15 and 30 frames, all told.

\section{Calculus Problems}

\subsection{Somewhat trickier problems}

\begin{frame}{Problem 11}
Evaluate
$$\lim_{n\rightarrow\infty}\left(\sqrt{1+2+3+\dots+n}-\sqrt{1+2+3+\dots+(n-1)}\right)$$
\end{frame}

\begin{frame}{Problem 12}
Evaluate
$$\lim_{n\rightarrow\infty}\frac{\sqrt{1\cdot 2} + \sqrt{2\cdot 3} + \dots + \sqrt{n(n+1)}}{n^2}$$
\end{frame}

\begin{frame}{Problem 13}
Evaluate
$$\sum_{n=1}^\infty\frac{1}{2^n}\sin\frac{n\pi}{3}.$$
\end{frame}

\begin{frame}{Problem 14}
Evaluate
$$\sum_{n=1}^\infty\frac{1}{2^n}\sin\frac{n\pi}{3}.$$
\end{frame}

\begin{frame}{Problem 15}
If $f(x) = e^{x^3} + e^{-x^3}$, evaluate $f^{(2022)}(0)$
\end{frame}

\begin{frame}{Problem 16}
Calculate
$$\lim_{n\rightarrow\infty}\sum_{i=1}^n\frac{1}{\sqrt{1^3+2^3+\dots+i^3}}$$
\end{frame}

\begin{frame}{Problem 17}
Calculate
$$\lim_{x\rightarrow\infty}\left(\sin\sqrt{x+\pi}-\sin\sqrt{x}\right)$$
\end{frame}

\begin{frame}{Problem 18}
Calculate
$$\lim_{n\rightarrow\infty}\cos\frac{1}{2}\cos\frac{1}{4}\cos\frac{1}{8}\dots\cos\frac{1}{2^n}$$
\end{frame}

\end{document}


